% !TEX root = ../main.tex
\section{The TokenCombine Operator}
\label{app:custom}

\textbf{Interface.} TokenCombine is a custom operator written against the SDK's C++ interface. Its inputs are a stage's
expert outputs in the layout of \S\ref{sec:red}, $[4,16C,d]$ in FP16, and per-SoC row indices $[4,Tk]$ that hold $-1$ for a
token's experts on other SoCs; its output is the per-SoC partials $[4,T,d]$. It thus replaces step 5's gather, mask and sum
over $k$, while the graph keeps the index construction, the add of the two stages' partials and the elementwise sum across
SoCs. Its configuration declares the output tileable along all three axes; for an output tile $[b_0{:}b_1,t_0{:}t_1,d_0{:}d_1]$
the operator requests the data tile $[b_0{:}b_1,0{:}16C,d_0{:}d_1]$ and the index tile $[b_0{:}b_1,t_0k{:}t_1k]$, so a tile
on SoC $c$ reads only $c$'s expert outputs. The kernel converts each indexed row from FP16 to FP32 with HVX vector
instructions, accumulates a token's rows in FP32 and converts the sum back to FP16 once; it asks for its operands in the
core's vector memory.

\textbf{Integration.} In a four-SoC test of one stage at $T=128$, the operator's output is bit-identical to step 5's
aggregation. In the full model at $T=512$ its tiles land on all 64 cores, and the cross-SoC traffic of the MoE layer stays
at 14.1\,MiB. A scalar version was 6.5--7.8$\times$ slower than step 5 on two layers, compute-bound at about 26\,ms of busy
time per core and layer; the HVX version runs the two layers in 14.8 against 15.7\,ms. The compiler did not place the data
input in vector memory and warns about it 26,112 times, without effect on the result. It instantiates the kernel at every
output tile and hands each SoC's code to one single-threaded \texttt{hexagon-clang++ -O3} run of about 2.7\,h, so a
full-model program takes 650--656 minutes to compile.

\textbf{Numerics.} The last-position logits differ from step 5's by $1.68\times10^{-2}$ relative $L_2$, with a largest
absolute difference of 0.273 for logits up to 24.7, the same top five and the same next token; both programs repeat
bit-identically, so the difference is deterministic. The kernel rounds each stage's FP32 partial to FP16 once. Step 5 adds
the two stages slot by slot, which is exact because each slot is non-zero in one stage only, and then reduces the eight
slots with the compiler's \texttt{ReduceSum}, whose accumulation precision is undocumented. Formulations that reorder the
same sums differ from step 5 by $0.7$--$3.3\times10^{-2}$; those that keep its order are bit-identical. One matching next
token is not evidence of general equivalence.
